\section*{Chapitre \ref{ch:s1:anatomy-of-a-stat-arb-book} --- Anatomie d'un portefeuille de stat arb}

\begin{solution}{exo:s1:anatomy-of-a-stat-arb-book:1}
Des paris indépendants s'ajoutent en variance~: $\sqrt{946 \times 252} \times 0.02 = 9.8$ positions équivalentes, à comparer avec $946 \times 252 = 238\,392$ jours-positions. Le
L'edge croît linéairement et le bruit avec la racine carrée, ce qui fait tout l'argument en faveur de la largeur (breadth).
\end{solution}

\begin{solution}{exo:s1:anatomy-of-a-stat-arb-book:2}
$0.04 - 0.0342 = 0.0058$~: 0.58~\% du capital par an, soit le prix du financement des positions longues au-delà du capital et des frais inclus dans le taux de rebate.
\end{solution}

\begin{solution}{exo:s1:anatomy-of-a-stat-arb-book:3}
$0.5 \times 1.5 + 0.5 \times 1.5 = 1.5$ fois le capital, que le portefeuille ne possède pas. Les capitaux propres doivent représenter la moitié de l'exposition brute, donc celle-ci vaut au plus $1/0.5 = 2.0$ fois
le capital~; une exigence de 15~\% de l'exposition brute en autorise $1/0.15 = 6.67$ fois.
\end{solution}

\begin{solution}{exo:s1:anatomy-of-a-stat-arb-book:4}
Le rebate vaut $1.5 \times (0.04 - 0.0025)$, le financement des longs $0.5 \times 0.045$~: $1.5 \times 0.0375 - 0.5 \times 0.045 = 0.0337$, soit 3.37~\% du capital par
an. Le portefeuille du chapitre rapporte 3.42~\% parce que la somme de ses longs plafonnés est un peu inférieure à 1.5 certains jours.
\end{solution}

\begin{solution}{exo:s1:anatomy-of-a-stat-arb-book:5}
L'exposition du signal aux facteurs de style change chaque jour (le momentum dérive, les tailles et les cours bougent), et la neutralisation quotidienne la retire à nouveau par des transactions. Le
portefeuille neutre au seul bêta et aux industries porte simplement l'exposition au style et ne la traite pas.
\end{solution}

\begin{solution}{exo:s1:anatomy-of-a-stat-arb-book:6}
Si $f_t$ est le coefficient de la régression en coupe et $e_t = r_t - X f_t$ son résidu, alors $r_t = X f_t + e_t$ par construction et
$w^\top r_t = w^\top X f_t + w^\top e_t = (X^\top w)^\top f_t + w^\top e_t$. L'égalité cesse de tenir lorsque les expositions utilisées dans l'attribution diffèrent de celles de la
régression (celles d'un autre jour ou d'un autre modèle), lorsque des titres absents de la régression sont détenus, ou lorsque le rendement du portefeuille est mesuré sur des
cours différents de ceux du modèle.
\end{solution}

\begin{solution}{exo:s1:anatomy-of-a-stat-arb-book:7}
Avec une demi-vie d'un jour, le portefeuille tourne de 129~\% de son capital par jour, verse 32.5~\% par an en coûts et rapporte 27.5~\% pour un ratio de Sharpe de 6.16~; avec cinq jours
56~\%, 14.0~\%, 11.7~\% et 2.47. Ici, aller plus vite est préférable, car le retournement simulé est fort et de courte durée et le coût de cinq points de base est faible~: ralentir
le portefeuille fait perdre plus de signal qu'il n'en économise en coûts. Avec des coûts réalistes, le classement peut s'inverser (chapitre~2).
\end{solution}

\begin{solution}{exo:s1:anatomy-of-a-stat-arb-book:8}
Un bêta nul et une exposition nulle en dollars laissent libres tous les autres facteurs. Le portefeuille du chapitre, neutre au bêta et aux industries, tirait 99.7~\% de sa variance des facteurs,
presque tous liés au momentum~; son rendement était un pari sur le style. La neutralité au marché n'est pas une neutralité aux facteurs, et l'alpha est ce qui subsiste après tous les facteurs que la firme peut acheter
à bas coût.
\end{solution}

\begin{solution}{pb:s1:anatomy-of-a-stat-arb-book:1}
\begin{enumerate}
\item Un mélange égal du retournement résiduel à un jour, normalisé en z-score, et du momentum 12--1 mois, lissé (demi-vie de deux jours), neutralisé par régression sur les
expositions du modèle de risque, mis à l'échelle pour une exposition brute de trois, plafonné à 1~\% par titre, sans vente à découvert sur les titres de la \omterm{def:s1:anatomy-of-a-stat-arb-book:htb}{liste des titres difficiles à emprunter} (hard-to-borrow).
\item 452 positions longues et 495 courtes en moyenne~; la position médiane représente 0.25~\% du capital, la plus importante 1.0~\%.
\item Le ratio d'information croît avec la racine carrée du nombre de paris indépendants~; la loi suppose des paris indépendants et ignore les coûts, et les expositions
communes du portefeuille et ses transactions quotidiennes violent ces deux hypothèses.
\item 1.9 million.
\item 23.6~\%, 23.6~\% et 1.00 neutre au bêta et aux industries~; 22.5~\%, 4.5~\% et 5.03 également neutre aux facteurs de style.
\item 99.7~\% pour le premier, presque tout du momentum (un P\&L de momentum de 9.0~\% par an avec une volatilité de 23.1~\%, et la valeur $-2.9~\%$)~; 7.6~\% pour le second.
\item Rotation de 91~\% du capital par jour contre 42~\%, et coûts de transaction de 22.9~\% par an contre 10.6~\%.
\item La simulation intègre un retournement fort et une dérive persistante à des coûts de cinq points de base~; les ratios de Sharpe publiés pour les stratégies de stat arb après coûts sont plus proches de 1.
\item Rebate 5.6~\%, financement des longs $-2.2~\%$, frais d'emprunt $-0.4~\%$ (inclus dans le rebate)~: net $+3.4~\%$ du capital par an.
\item Des capitaux propres égaux à la moitié de l'exposition brute, donc au plus le double du capital~; une exposition brute de trois exige une marge de portefeuille fondée sur le risque (un choc de 15~\% sans compensation exige
0.45 fois le capital).
\item Ces titres sont signalés avant l'optimisation et ne sont pas vendus à découvert~; s'ils l'avaient été, leurs frais (médiane voisine de 3~\%) auraient réduit le rebate et exposé le portefeuille au rappel.
\item \textbf{Résultat nommé.} 99.7~\% de la variance du P\&L du portefeuille non neutre aux facteurs de style provenait des expositions aux facteurs~; la charge du financement (financement des longs et
frais d'emprunt) représentait 6.2~\% de l'alpha brut de 42.0~\% du portefeuille neutre.
\item Ratios de Sharpe après coûts de 1.44 sur 1997--2007 pour leurs stratégies fondées sur l'ACP, 0.9 sur 2003--2007~; la version fonds indiciel coté (ETF) s'est dégradée de façon similaire après 2002.
\item Les portefeuilles construits de cette façon détenaient des positions similaires et ont perdu ensemble lorsque certains d'entre eux ont été dénoués.
\item Données vérifiées et stockées, modèle de risque mis à jour, signaux mélangés et lissés, portefeuille optimisé, transactions exécutées à la séance suivante, financement provisionné, le P\&L
attribué et réconcilié.
\item Parce qu'une partie qui ne se réconcilie pas signale une position, un cours, une exposition ou une charge erronés~: un résidu inexpliqué est un défaut, et non une ligne du rapport.
\item Des coûts réalistes incluant l'impact (livre~7, chapitre~27), les locates et les rappels, un vrai modèle de risque, et des signaux dont la force est estimée et non imposée.
\item Attribuer son P\&L~: une part importante de la variance venant des rendements des facteurs, ou une forte corrélation avec des indices de style, signale un pari sur un facteur.
\item Quelle part de sa variance provient de facteurs qu'il ne prétend pas détenir.
\item Un portefeuille long--short composé de nombreux petits paris, individuellement peu fiables, neutre à l'égard de ce qu'il ne veut pas détenir, dont le rendement vient de leur nombre.
\end{enumerate}
\end{solution}

\begin{solution}{iq:s1:anatomy-of-a-stat-arb-book:1}
Une exposition nette en dollars nulle et un bêta nul, donc le mouvement du marché ne le déplace pas. Il reste exposé aux industries, aux facteurs de style (momentum, valeur, taille), à l'encombrement (crowding) lié à des portefeuilles
similaires, ainsi qu'aux conditions de financement et d'emprunt.
\end{solution}

\begin{solution}{iq:s1:anatomy-of-a-stat-arb-book:2}
Les longs sont financés par du capital et par un solde débiteur auprès du prime broker~; les ventes à découvert portent sur des titres empruntés, le produit de la vente étant déposé en collatéral et rémunéré au taux de l'indice de référence diminué des
frais d'emprunt. Ce qui peut mal tourner~: les frais montent ou des titres sont rappelés, les appels de marge s'accroissent en période de stress, le courtier réduit son financement et le portefeuille est contraint de se réduire au pire
moment.
\end{solution}

\begin{solution}{iq:s1:anatomy-of-a-stat-arb-book:3}
Le ratio d'information vaut approximativement le coefficient d'information (IC) multiplié par la racine carrée de la largeur (breadth). Un faible pouvoir prédictif appliqué à de nombreux
titres indépendants et rééquilibrés souvent donne un ratio élevé, ce qui constitue la prémisse de la stat arb~; ses limites tiennent à la corrélation des paris et aux coûts de transaction.
\end{solution}

\begin{solution}{iq:s1:anatomy-of-a-stat-arb-book:4}
Davantage de capital investi dans les mêmes signaux (l'edge se partage et les coûts montent), des fourchettes plus étroites et des marchés plus rapides qui ont fait disparaître certains écarts de prix, et un encombrement (crowding) qui
a rendu les portefeuilles assez semblables pour perdre ensemble.
\end{solution}

\begin{solution}{iq:s1:anatomy-of-a-stat-arb-book:5}
Positions et cours tels qu'ils étaient alors connus, expositions du modèle de risque et rendements des facteurs~; P\&L décomposé en facteur (par facteur), spécifique, coûts (par transaction) et financement (par
titre). Il doit se réconcilier avec les positions multipliées par les rendements et avec le relevé de trésorerie du prime broker, portant sur les frais et les rebates.
\end{solution}

\begin{solution}{iq:s1:anatomy-of-a-stat-arb-book:6}
Chaque position vaut $G/N$~; la variance spécifique est $N (G/N)^2 \sigma^2$, donc la volatilité vaut $G\sigma/\sqrt N$. Le rendement espéré est $N (G/N)\alpha = G\alpha$,
et le ratio vaut $\alpha\sqrt N/\sigma$~: il croît avec $\sqrt N$ à edge par titre constant, indépendamment de $G$.
\end{solution}
